\subsection{Conjunctions and the literature on existential rank}
\label{subsec:intersection-rank}

The general relation between witness number and finite fibers is part of
the theory of existential rank and essential fiber dimension developed
by Daans, Dittmann, and Fehm \cite[Corollary 4.21]{DDF2025}. Here
\emph{existential rank} means the minimum number of existential
quantifiers in a definition of a specified set. Our finite-fiber proof
is a direct constructive treatment of the special case needed here,
including exceptional strata and nonzero conditions. We make no claim
that the general quantifier-saving phenomenon is new.

For clarity, it gives the following useful known intersection bound
\cite[Theorem 1.4]{DDF2025}, also stated in
\cite[Theorem 2.5]{Daans2024}, with an explicit constructive proof over
the rationals.

\begin{proposition}\label{prop:intersection-rank}
Let $D_i\subseteq\Q^n$, for $1\le i\le r$ and $r\ge1$, be
defined by one polynomial equation with $m_i\ge1$ existential
witnesses. Their intersection has a one-polynomial definition with at
most
\[
 1+\sum_{i=1}^r(m_i-1)
\]
witnesses. Such a definition can be constructed effectively from the
given polynomials.
\end{proposition}
\begin{proof}
Write the $i$th polynomial as
\[
 F_i(\mathbf x,\mathbf u_i,Z_i)
   =\sum_{j=0}^{d_i}a_{ij}(\mathbf x,\mathbf u_i)Z_i^j,
\]
where $\mathbf u_i$ has $m_i-1$ coordinates and the witness lists
are private to their respective predicates. Let $C_i$ be the closed
condition that all coefficients $a_{ij}$ vanish. Replace the zero set
of $F_i$ by the constructible set
\[
 (\{F_i=0\}\setminus C_i)\ \cup\ (C_i\cap\{Z_i=0\}).
\]
This replacement preserves the projection to
$(\mathbf x,\mathbf u_i)$ over any field: a nonzero polynomial has
its usual roots, and an identically zero polynomial permits the choice
$Z_i=0$. Each geometric fiber in $Z_i$ is now finite. Nonzero
constant specializations give empty fibers and cause no exception.

Take the fiber product of these presentations over $\mathbf x$.
It has finite geometric fibers over
$(\mathbf x,\mathbf u_1,\ldots,\mathbf u_r)$.
Theorem~\ref{thm:finite-fibre-one-variable} replaces all the $Z_i$
by one additional variable. Quantifying the remaining
$\sum_i(m_i-1)$ private coordinates proves the formula and its
effectiveness. Genuinely shared coordinates remain in $\mathbf x$;
they are never silently duplicated or identified.
\end{proof}

For example, each of the forty-two pole predicates has a four-witness
presentation after retaining its three free pole parameters and encoding
its remaining finite choices by one variable. The proposition combines
them with $42\cdot4-41=127$ witnesses. This is exactly the $126$
retained pole coordinates and one common finite-choice variable already
used in the global finite-fiber construction. It does not provide another
forty-one-variable saving on top of that construction.
